\section{Materials and methods}\label{sec:methods}

\subsection{Global cropland sites, weather and crop calendars}\label{sec:data}
We studied five staple crops that together cover about 670\,Mha of harvested area: wheat, rice, maize, soybean and potato. Harvested areas (rainfed and irrigated) were taken from the SPAM 2010 gridded production maps \citep{yu2020} and aggregated to 1\textdegree{} cells. For each crop, 80 cells with at least 5000\,ha of the crop were drawn without replacement with probability proportional to harvested area. An unweighted average over the sample therefore approximates an area-weighted average over global production. The 400 crop--site pairs lie in 381 distinct cells in 50 countries (Fig.~\ref{fig:map}). Sowing and maturity dates came from the GGCMI Phase~3 crop calendar at 0.5\textdegree{} \citep{jagermeyr2021}. For wheat we used the winter-wheat calendar unless the coldest-month mean temperature was below $-8$\,\textdegree C, in which case the spring-wheat calendar was used. Daily weather for 2001--2020 was retrieved from NASA POWER \citep{sparks2018}: global and diffuse shortwave irradiance, maximum, minimum and dew-point temperature, precipitation and wind speed at 2\,m. Each site contributes 19 seasons (sowing years 2001--2019), each simulated as the 365 days after sowing, which gives 7600 site-seasons.

Climate change was represented by the global warming levels +1.5, +2 and +3\,\textdegree C relative to 2001--2020. Daily maximum, minimum and dew-point temperatures were shifted uniformly, which keeps relative humidity approximately constant and raises the vapour-pressure deficit. The atmospheric CO$_2$ concentration was set to 390, 450, 500 and 600\,ppm for the baseline and the three warming levels. Precipitation and irradiance were left unchanged. This delta approach isolates the thermal signal of warming; its limitations are discussed in Section~\ref{sec:limits}.

\subsection{A differentiable agrivoltaic simulator}\label{sec:sim}
The simulator advances all site-seasons in parallel with a daily time step for the crop and an hourly time step for radiation and PV (Fig.~\ref{fig:scheme}). All equations are implemented as PyTorch tensor operations \citep{paszke2019}, so that yields and electricity are differentiable with respect to the design and control variables.

\paragraph{Hourly irradiance}
Daily global and diffuse irradiance were disaggregated into hourly values with the \citet{collares1979} and \citet{liu1960} ratios and renormalised to the daily totals. Beam irradiance on the horizontal is the difference between the two. Solar position follows the standard declination and hour-angle equations.

\paragraph{Array geometry and PV electricity}
We consider elevated horizontal single-axis trackers with north--south axes. The axis height is $1.25$ module widths, which leaves space for machinery under the modules. The rotation $\phi$ is limited to $\pm60^\circ$, and the ground coverage ratio $g$ is the module width divided by the row pitch. Let $\psi$ be the solar zenith angle projected onto the east--west vertical plane. The energy-optimal rotation with backtracking \citep{lorenzo2011} is $\phi_{\mathrm{bt}}=\psi-\mathrm{sgn}(\psi)\arccos(\cos\psi/g)$ when $\cos\psi<g$, and $\phi_{\mathrm{bt}}=\psi$ otherwise. The \emph{light-sharing} rotation $\phi_{\mathrm{sh}}$ is the feasible angle closest to edge-on to the sun, which minimises the shadow on the crop. Given a light-sharing level $u\in[0,1]$ chosen once per day, the rotation in every hour is
\begin{equation}
\phi = (1-u)\,\phi_{\mathrm{bt}} + u\,\phi_{\mathrm{sh}}.
\label{eq:phi}
\end{equation}
Plane-of-array irradiance combines beam irradiance on the module (reduced by the fraction of the module shaded by the neighbouring row, $\max\{0,1-\cos\psi/(g\,|\cos(\phi-\psi)|)\}$), isotropic sky diffuse and ground-reflected light (albedo 0.2). The cell temperature follows the NOCT model (45\,\textdegree C), with an hourly air temperature interpolated sinusoidally between the daily extremes. Electricity per unit land is plane-of-array irradiance $\times\,\eta\times g$, with $\eta=0.21\,[1-0.0035(T_{\mathrm{cell}}-25)]\,(1-0.14)$ for module efficiency, temperature losses and system losses. A conventional ground-mounted PV plant ($g=0.4$, backtracking, no crop) is the energy reference $E_{\mathrm{PV}}$.

\paragraph{Light and microclimate at the crop}
Averaged over the row pitch, the fraction of ground reached by beam light is $1-\min\{1,\,g\,|\cos(\phi-\psi)|/\cos\psi\}$. The fraction of sky diffuse reaching the ground is a two-dimensional isotropic sky-view factor. We pre-computed it for 25 values of $g$ and 13 rotation angles by ray casting over 31 neighbouring rows, and interpolate it bilinearly (differentiably). The daily ratio $s$ of the irradiance at the crop to that of the open field defines the crop's shading. Following observations of cooler canopies under PV modules \citep{marrou2013,barron2019}, the canopy maximum temperature is lowered by $\kappa(1-s)$ with $\kappa=2$\,\textdegree C (a sensitivity analysis over $\kappa$ is reported in Section~\ref{sec:sensitivity}). The reference evapotranspiration under the modules, ET$_0$, is recomputed with the FAO-56 Penman--Monteith equation \citep{allen1998} from the reduced radiation and canopy temperature. Soil water conservation under the modules \citep{adeh2018,elamri2018} therefore emerges from the model and is not imposed.

\paragraph{Crop growth}
Crop growth follows the SIMPLE model \citep{zhao2019} with the species parameters of its Table~1a (Table~\ref{tab:params}). Daily biomass growth is
\begin{equation}
\Delta B = R_{\mathrm{crop}}\cdot \mathrm{RUE}\cdot f_{\mathrm{Solar}}\cdot f_{\mathrm{CO_2}}\cdot f_{\mathrm{Temp}}\cdot \min\{f_{\mathrm{Heat}},f_{\mathrm{Water}}\},
\label{eq:simple}
\end{equation}
where $R_{\mathrm{crop}}=s\,R$ is the radiation received by the crop. $f_{\mathrm{Solar}}$ is the canopy's radiation interception, driven by thermal time and accelerated senescence under heat and drought stress. $f_{\mathrm{Heat}}$ decreases linearly between the heat threshold $T_{\mathrm{heat}}$ and the extreme temperature $T_{\mathrm{ext}}$ of the canopy maximum temperature. $f_{\mathrm{Water}}=1-S_{\mathrm{water}}\,\mathrm{ARID}$, with the agricultural reference index for drought \citep{woli2012} computed from a daily soil-water balance with curve-number runoff, drainage and root water uptake. Irrigated fields have $f_{\mathrm{Water}}=1$. Yields of the rainfed and irrigated parts of a site are combined with the SPAM irrigated fraction. The thermal-time requirement to maturity was calibrated for every site as the mean thermal time between the GGCMI sowing and maturity dates under the 2001--2020 climate, and the canopy parameters were scaled accordingly. The cultivar is thus adapted to the present climate, and warming shortens the season. Growth stops smoothly once the thermal-time requirement is met, and yield is biomass multiplied by the harvest index. Because SIMPLE uses a linear radiation-use efficiency, it does not credit the higher light-use efficiency of shaded canopies. Shading losses are therefore, if anything, overestimated.

\subsection{SALS: stress-aware light sharing}\label{sec:sals}
SALS learns two functions (Algorithm~\ref{alg:sals}):
\begin{itemize}
\item a \emph{design network} $g=D_\theta(\mathbf{z})$, a multilayer perceptron with two hidden layers of 32 units, which maps site descriptors $\mathbf{z}$ to a ground coverage ratio in $[0.05,0.60]$. The descriptors are the crop (one-hot), latitude, mean irradiance, growing-season maximum temperature, precipitation and vapour-pressure deficit, the frequency of heat-stress days and the irrigated fraction;
\item a \emph{control network} $u_t=\pi_\omega(\mathbf{x}_t)$, a multilayer perceptron with two hidden layers of 64 units and a sigmoid output, which chooses the daily light-sharing level of Eq.~\eqref{eq:phi} from information available on day $t$: relative thermal time $\mathrm{TT}/T_{\mathrm{sum}}$, radiation interception $f_{\mathrm{Solar}}$, relative soil water, the previous day's ARID, the day's forecast of maximum temperature relative to $T_{\mathrm{heat}}$ and $T_{\mathrm{opt}}$, irradiance and ET$_0$, the irrigation status, the design $g$, latitude, day of season and the crop.
\end{itemize}
Outside the cropping season the trackers backtrack for maximum electricity. For site-season $i$ with open-field yield $Y^{\mathrm{open}}_i$, let $r_i=Y_i/Y^{\mathrm{open}}_i$ be the yield retention and $e_i=E_i/E_{\mathrm{PV},i}$ the electricity relative to a conventional PV plant. SALS solves
\begin{equation}
\max_{\theta,\omega}\;\mathbb{E}_i\big[e_i\big]\quad\text{s.t.}\quad r_i\ge\rho\;\;\forall i,
\label{eq:problem}
\end{equation}
where $\rho$ is the food-security floor ($\rho=0.9$ unless stated otherwise). It does so with the penalty objective
\begin{equation}
\mathcal{L}(\theta,\omega)=\mathbb{E}_{i\sim\mathcal{B}}\Big[-e_i+\mu\,\max\{0,\rho-r_i\}^2\Big],\qquad \mu=30.
\label{eq:loss}
\end{equation}
Because $r_i$ and $e_i$ are outputs of the differentiable simulator, the exact gradient of Eq.~\eqref{eq:loss} with respect to $(\theta,\omega)$ is obtained by back-propagation through the 365 daily steps of every season (back-propagation through time). This is the key difference from reinforcement learning with black-box simulators \citep{gautron2022}, which estimates gradients from sampled returns. Each minibatch contains 512 site-seasons, each duplicated into a rainfed and an irrigated field, drawn from one of the four climates in turn. The networks are trained with Adam \citep{kingma2015} (learning rate $3\times10^{-3}$, cosine decay, gradient-norm clipping at 1) for 800 steps. The rare updates with a non-finite gradient are skipped.

\begin{algorithm}[t]
\caption{SALS training}\label{alg:sals}
\begin{algorithmic}[1]
\Require training site-seasons in four climates; floor $\rho$; penalty $\mu$
\State initialise design network $D_\theta$ and control network $\pi_\omega$
\For{step $=1,\dots,N$}
  \State sample a climate and a minibatch $\mathcal{B}$ of site-seasons (rainfed and irrigated rows)
  \State $g_i\gets D_\theta(\mathbf{z}_i)$; initialise crop and soil states
  \For{day $t=1,\dots,365$}
    \State $u_{it}\gets\pi_\omega(\mathbf{x}_{it})$ if the crop is in the field, else $0$
    \State hourly rotation (Eq.~\ref{eq:phi}), PV electricity, crop light, canopy temperature, ET$_0$
    \State advance the soil-water balance and the SIMPLE crop state (Eq.~\ref{eq:simple})
  \EndFor
  \State $r_i\gets Y_i/Y_i^{\mathrm{open}}$, $e_i\gets E_i/E_{\mathrm{PV},i}$; loss $\mathcal{L}$ (Eq.~\ref{eq:loss})
  \State back-propagate through the season; Adam update of $(\theta,\omega)$
\EndFor
\end{algorithmic}
\end{algorithm}

\subsection{Baselines}\label{sec:baselines}
We compared SALS with the following references and baselines.
\begin{itemize}
\item \emph{Open field}: no modules ($r=1$, $e=0$, LER $=1$).
\item \emph{PV plant}: conventional ground-mounted PV without a crop ($e=1$, LER $=1$).
\item \emph{AV static (rule)}: backtracking trackers with one GCR per crop, the largest value on a grid $\{0.05,0.10,\dots,0.50\}$ whose mean retention on the training seasons reaches $\rho$. This mimics a design guideline.
\item \emph{AV static (oracle GCR)}: backtracking trackers with the largest GCR that meets $\rho$ on average over each test site's own seasons. The GCR is chosen with knowledge of the test weather.
\item \emph{Seasonal sharing}: the light-sharing rotation ($u=1$) throughout the cropping season, with the per-site oracle GCR.
\item \emph{Stress rule}: an expert rule that shares light ($u=1$) during the main growth phase ($0.1<\mathrm{TT}/T_{\mathrm{sum}}<0.9$) except on days with a heat-stress forecast ($T_{\max}\ge T_{\mathrm{heat}}$) or drought (ARID\,$>0.5$), when the modules track the sun and shade the crop. It uses the per-site oracle GCR.
\item \emph{Oracle open-loop}: the GCR and all 365 daily controls of every test season are optimised directly by gradient descent (150 Adam steps) with perfect knowledge of that season's weather. This is an upper bound that no deployable controller can reach.
\end{itemize}
The three oracle-GCR baselines therefore have access to the test weather, whereas SALS chooses its design from climate descriptors alone.

\subsection{Evaluation protocol and metrics}\label{sec:protocol}
Sites were split spatially: 10\textdegree{}$\times$10\textdegree{} blocks were randomly assigned to training (70\%) or test (30\%), so that no test site has a training site in the same block. Seasons sown in 2001--2014 at training sites were used for training, and seasons sown in 2015--2019 at test sites for evaluation, so that the test set is out of sample in both space and time. We report the mean electricity (MWh\,ha$^{-1}$\,yr$^{-1}$) and relative electricity $e$, the mean yield retention $r$, the \emph{compliance} (share of test seasons with $r\ge\rho$), and the land-equivalent ratio LER $=r+e$ \citep{dupraz2011}. Differences between methods were tested with paired Wilcoxon signed-rank tests over test seasons, and uncertainty is given as bootstrap 95\% confidence intervals over sites.
